\section{A finite menu of plugs and dense pins}\label{sec:pins}

We construct a finite menu of local endpoint models that occurs densely
along every relevant macroscopic rim: each sufficiently large subarc
contains a usable endpoint from that menu. These are the models used in
Proposition~\ref{prop:endpoint-potentials}.  Its geometric constants and the
thresholds $\beta,\beta'>0$ will be independent of the plug scale $h$.
Placements are translates by $h\Z^2$, in one common frame, of the scaled
members of this menu.  Both color choices are included.  Throughout this
section a probability assertion is interpreted with $\eps\to0$ first, at
fixed $h$, and subsequently $h\to0$.  Put
\[
                    r_h=h^{1-\kappa},\qquad 0<\kappa<1.
\]
We will choose $\kappa$ sufficiently small in
Section~\ref{sec:actual-arcs}.

\subsection{The qualifications at a pin}

\begin{definition}[Qualified pin]\label{pin:qualification}
Let $P$ be an oriented rim in a bulk patch.  A qualified pin consists of a
vertex $p\in P$, a grid center $z$ with $|p-z|\le h$, a placed menu type,
its root ranks, and the following data.  We describe the black case;
interchanging colors gives the white case.
\begin{enumerate}[label=(\roman*)]
\item In the original configuration there are three surrounding gate
circuits, in prescribed annuli $[a_kh,b_kh]$, $k=1,2,3$.  The first two
are the annuli of the plug.  The third has $a_3>100b_2$ and isolates an
ordered central subarc $\Delta$ through $p$: all visits of $P$ to
$B(z,a_3h)$ belong to $\Delta$, and $\diam\Delta\le2b_3h$.
\item In the nominal chart the local ordered traversal of $P$ follows
the horizontal midline from before $u=-8$ to after $u=8$, with transverse
error and longitudinal backtracking less than $0.1$.  The arc $\Delta$
lies in $|u|<1$, and its marked point is near $u=0$.  There are no other
visits to $H$.  The fit test has the closed-limit property of
Lemma~\ref{lemma:exact-fit} below.
\item A prescribed local preparation, applied to the original sample
using an auxiliary seed attached to this chart, placement, and color,
only recolors existing sites in the opposing strip from black to
white, keeping the point configuration fixed.  It leaves the
fitted black traversal unchanged.  Its inner state $\sigma'$ belongs
to $G$, and the chosen type is admitted.  The black root contains vertices of the
\emph{actual traversal} near $u=-1.7$ and $u=1.7$, with error less than
$0.05$, joined along that traversal in the root window.  The two black
side passages are also supplied by that traversal.  All these paths are
retained paths.
\item For the white color separately, each of the two ordinary side-test
probabilities, conditional on $\sigma'$ \emph{and the actual point
geometry}, is at least $\beta$.
\end{enumerate}
The prepared state is shared whenever the same candidate choice is used;
preparations for other choices start again from the original sample.
\end{definition}

The last qualification is stronger than the annealed propensity condition
in $G$.  It will permit an exterior white forcing while retaining the
actual point geometry.  We will not need a corresponding quenched lower
bound for the black side passages, since those passages are already
present on $P$.

\subsection{Three gates and local inverse-diameter tests}

Work in units of $h$, centered at one grid point.  Choose integers $n$
and a large constant $M$, and set
\[
 A_1=30,\qquad B_i=MA_i,\qquad A_{i+1}=10^4B_i,
 \qquad L=B_n.
\]
In each trial annulus take slightly inset primal boundary cycles and
test for four disjoint arms with both colors represented.  The tests use
retained graphs in separated annular collars; local density or mesh
failure is included as a positive test.  Thus different trials use
independent restrictions of the marked Poisson process.  By
Lemma~\ref{lemma:mixed-four}, each test has limiting probability at most
$CM^{-1-c}$, for some $c>0$.

Suppose a rim has a vertex within distance one of the center and extends
well beyond the trials.  Besides its own color arm it supplies an
opposing arm: the facing neighboring vertices near the mark and beyond
the trial belong to one opposing cluster by Lemma~\ref{lemma:rim}.
Trim its connection to the trial annulus.  This is a local arm even when
the rim is completed on the sphere.  If there are three disjoint black
crossings of a trial annulus, this opposing arm gives the positive
four-arm test.  A negative test therefore permits the two-gate
construction of Lemma~\ref{lemma:two-gates}.

Call a trial clean when its test is negative.  The probability that fewer
than three trials are clean, allowing either rim color, is bounded in
the microscopic limit by
\begin{equation}\label{pin:trials}
                   Cn^2(CM^{-1-c})^{n-2}.
\end{equation}
Indeed at least $n-2$ prescribed trials must then be positive, and there
are at most $Cn^2$ ways to choose the remaining indices.  From three
clean trials choose gate circuits in their increasing order.  Denote the
corresponding inset annular radii by $a_k,b_k$.  The large separation
allows a cutoff $m$ satisfying all the inequalities in
Section~\ref{sec:plugs}; it also gives $a_3>100b_2$.

A simple black rim visiting the inside and outside of the third circuit
meets it at its two black vertices, once at each and nowhere else.  Its
inner arc $\Delta$ contains the marked vertex, is contained in
$\overline B(0,b_3)$, and contains every visit to $B(0,a_3)$.  This
description remains true after adding whites away from the traversal:
the circuit still has at most two black vertices, and its two contacts
with the unchanged traversal remain black.

Extend $\Delta$ in both directions to the nearest hits of
$\partial B(0,10L)$ and call the resulting ordered arc $K$.  Mark its
two endpoints, the two endpoints of $\Delta$, and the vertex $p$.  The
complementary rim travels beyond radius $100L$ for all sufficiently small
$h$ in the application.  Let $E$ be its portion in
$\overline B(0,50L)$, together with $\partial B(0,50L)$.  Include the two
endpoints of $K$ in $E$.  This compact set is connected: each component
of the truncated complementary path meets the added circle, or is joined
to it along that path.  A \emph{profile} is $(K,E)$ with these marks and
the finite trial-index data.  We use marked Fr\'echet distance for $K$
and Hausdorff distance for $E$.  A Fr\'echet comparison allows an
increasing reparametrization matching the marks.

Figure~\ref{fig:pin-profile} separates the ordered arc from the
complementary data.  Adding the outer circle keeps $E$ connected, while
the third gate isolates the middle interval $\Delta$ from every other
visit to the inner region.  The inverse-diameter tests below will
preserve this separation when profiles converge.

\begin{figure}[tb]
\centering
\begin{tikzpicture}[scale=.95,>=Stealth,font=\small]
 \draw[violet!75!black,thick] (0,0) circle (2.4);
 \draw[gray,dashed] (0,0) circle (1.45);
 \draw[orange!75!black,thick] (0,0) circle (.55);
 \draw[very thick] (-1.45,0)--(-1.18,.17)--(-.93,-.12)
       --(-.72,.08)--(-.55,0);
 \draw[very thick,blue!65!black] (-.55,0)--(-.3,-.15)
       --(0,0)--(.3,.15)--(.55,0);
 \draw[very thick,->] (.55,0)--(.79,-.1)--(1.12,.18)--(1.45,0);
 \draw[violet!75!black,very thick] (-1.45,0)--(-1.72,.17)
       --(-2.04,-.13)--(-2.4,0);
 \draw[violet!75!black,very thick] (1.45,0)--(1.75,-.16)
       --(2.06,.12)--(2.4,0);
 \fill[violet!75!black] (-1.45,0) circle (2.3pt)
       (1.45,0) circle (2.3pt);
 \fill[black] (0,0) circle (1.8pt);
 \fill[blue!65!black] (-.55,0) circle (1.8pt)
       (.55,0) circle (1.8pt);
 \node[above left] at (0,.03) {$p$};
 \node[blue!65!black] at (0,-.85) {$\Delta$};
 \draw[->,blue!65!black] (0,-.65)--(-.25,-.16);
 \node at (1.03,.56) {$K$};
 \node[orange!75!black,anchor=west] at (2.65,.85) {third gate};
 \draw[->,gray] (2.58,.73)--(1.0,.92)--(.36,.42);
 \node[gray,anchor=east] at (-2.65,1.15) {$\partial B(0,10L)$};
 \draw[->,gray] (-2.58,1.05)--(-1.03,1.03);
 \node[violet!75!black,anchor=west] at (1.85,1.85)
       {$\partial B(0,50L)\subset E$};
 \draw[->,violet!75!black] (1.83,1.72)--(1.69,1.69);
 \node[violet!75!black,align=center] at (0,-3.05)
       {the two terminal points of $K$\\also belong to $E$};
 \draw[->,gray] (-1.04,-2.78)--(-1.44,-.08);
 \draw[->,gray] (1.04,-2.78)--(1.44,-.08);
\end{tikzpicture}
\caption{A schematic profile in units of $h$, with radial scales not to
scale.  The blue interval $\Delta$ is the part of $K$ inside the third
gate.  The purple set $E$ consists of the complementary rim inside the
outer circle, including the two endpoints of $K$, and the entire outer
circle.  The omitted complementary rim reaches beyond radius $100L$.
Additional excursions are allowed; no smoothness of a profile is assumed.}
\label{fig:pin-profile}
\end{figure}

Fix $\eta_j\downarrow0$, with $\eta_1\le1$.  Choose $\delta_j>0$
successively so small that the following forbidden-return tests have a
summably small total probability.  Two visits in
$\overline B(0,60L)$ at distance at most $\delta_j$ must not have
departures to distance $\eta_j$ along \emph{both} intervening directions
of the rim.  To flag violations use only the local necessary arm events
in Lemma~\ref{lemma:near-return}, on a $\delta_j$-grid inside
$B(0,100L)$, together with the relevant retained-mesh tests.  Their
limiting total probability at level $j$ is at most
\begin{equation}\label{pin:inverse-cost}
 C(L/\delta_j)^2(\delta_j/\eta_j)^{2+c}
       =CL^2\delta_j^c\eta_j^{-2-c}.
\end{equation}
It can be made arbitrarily small.  The events, unlike the profiles
themselves, depend only on this bounded local Poisson restriction.
Only a finite initial list of levels will ever be imposed at a fixed
$h$.

\begin{lemma}[Compactness of profiles]\label{lemma:profile-compactness}
Every sequence of the above profiles satisfying increasingly long
initial lists of the inverse-diameter constraints has a subsequence
converging in marked Fr\'echet and Hausdorff distance to a simple marked
arc $K$ and a compact connected set $E$.  The required marks remain
distinct.  In the limit all visits of $K$ to
$\overline B(0,4b_2)$ occur in the interior of its marked middle interval
$\Delta\subset\overline B(0,b_3)$.  The set $E$ misses
$\overline B(0,4b_2)$ and misses $K\cap B(0,6L)$.
The family of all such limiting profiles is compact.
\end{lemma}

\begin{proof}
The complementary portion of the rim makes a fixed large departure.
Consequently, if two points of $K$ are sufficiently close, the portion
between them in $K$ has small diameter; otherwise both directions would
violate one of the imposed constraints.  More precisely, for a fixed
$j$, closeness at scale $\delta_j$ forces this portion to have diameter
at most $4\eta_j$, after requiring $\delta_j\le\eta_j$.  For tests
initially formulated at vertices, use the buffered test and mesh
smaller than a fixed fraction of each of the finitely many used
$\delta_j$'s to obtain the same implication for points on the polygonal
arc.  The harmless factor four can be absorbed by starting with a
correspondingly smaller sequence $\eta_j$.

Here is a finite-net argument that does not assume length bounds.  Fix
a displacement scale $d>4\eta_j$.  Starting at the first endpoint, stop
successively when the curve first reaches distance $d$ from the preceding
stop, also inserting the marked times.  Apart from a bounded number of
stops adjacent to inserted marks, different stop starts are
$\delta_j$-separated: an intervening displacement $d$ rules out their
closeness by the preceding inverse-diameter property.  Packing in the
fixed disk bounds the number of stops.  Every portion between stops is
contained in a $d$-ball about its start.  Joining the stops by segments
and rounding their positions to a sufficiently fine finite grid gives
a finite marked Fr\'echet net, with error tending to zero as $d\to0$.

Choose a rapidly Cauchy subsequence in these nets and align successive
parametrizations by increasing homeomorphisms matching the marked
times.  The aligned curves converge uniformly to a continuous curve.
If two of its times have the same image, the intervening portion has
arbitrarily small diameter by the inverse constraints, and hence is
constant.  Every point fiber is thus an interval.  Collapse these
constant intervals.  The quotient is a nontrivial compact connected
metrizable linear order, with its order topology: order rays have open
preimages, and compactness identifies the quotient topology.  Such an
order is an interval, as is seen by embedding a countable dense suborder
in an interval and completing the order.  The induced parametrization
is injective and continuous, so the resulting arc is simple.  A
nondecreasing parametrization with constant intervals can be uniformly
approximated by increasing homeomorphisms; with finitely many distinct
marks this can be done separately between consecutive marks.  More
explicitly, after choosing the quotient parametrization so that its
nondecreasing map $q$ fixes the normalized mark times, the maps
$(1-t)q+t\,\mathrm{id}$, $0<t<1$, are increasing homeomorphisms fixing
those times and converge uniformly to $q$.  Thus the
convergence is in the asserted marked Fr\'echet topology.

The radial requirements keep the marks distinct.  For example, if the
ends of $\Delta$ approached one another, the path through the core
would make a fixed departure in one direction and the complementary
path would do so in the other, contradicting an inverse constraint.
The same argument treats the endpoints of $K$ and every pair of
nonadjacent marks; adjacent marks have the radial separation furnished
by the core and the two circles.  In particular the limit is not
constant.

Pass also to a Hausdorff subsequence for $E$.  Connectedness is preserved
under Hausdorff limits of nonempty compact connected sets.  Before the
limit, all visits to $B(0,a_3)$ belong to $\Delta$, and
$a_3>100b_2$; this gives the assertion about
$\overline B(0,4b_2)$ with a margin.  If a complementary-rim point
approached a point of $K$ in $B(0,6L)$, both intervening rim portions
would pass through a terminal hit at radius $10L$.  A fixed
inverse-diameter constraint excludes this approach.  It excludes a
limiting intersection as well.  The added circle at radius $50L$ is
irrelevant in these inner regions.

Finally take any sequence of limiting profiles.  Approximate its $j$th
member closely by a prelimit profile satisfying at least the first $j$
constraints, with the same finite trial data after passing to a
subsequence.  The argument just proved gives a convergent subsequence
of these approximants, and therefore of the limiting profiles.  This
proves compactness of the limiting family.
\end{proof}

\subsection{Models that fit exactly after a small ambient change}

\begin{lemma}[Exact fitting at closed limits]\label{lemma:exact-fit}
Each limiting profile in Lemma~\ref{lemma:profile-compactness} has a
nominal plug model, fixed local stubs, a positive allowed chart-change
tolerance, and an open fit neighborhood with the following properties.
Every fitting prelimit traversal has the strict tracking properties in
Definition~\ref{pin:qualification}.  For every simple profile in the
closure of this neighborhood there is an orientation-preserving ambient
homeomorphism $g$, supported in $B(0,50L)$ and within the allowed
tolerance, such that its ordered arc agrees \emph{exactly} with
$g\Phi(u,0)$ through $gH$ and continues along it to $|u|=9.5$.
There are no other visits of the arc or its unwanted set $E$ to $gH$.
A finite collection of such neighborhoods contains all profiles
satisfying a sufficiently long finite initial list of constraints.
\end{lemma}

\begin{proof}
Extend $K$ to a proper line by outward rays at its two distinct endpoints
on $\partial B(0,10L)$.  The rays are disjoint and meet the closed disk
only at those endpoints.  In the one-point compactification of the
plane, this proper line together with infinity is a Jordan curve on
the sphere.  Apply the sphere Jordan--Schoenflies extension theorem
\cite[Theorems~10.2--10.3, p.~71]{Moise1977} to a correspondence from
the compactified horizontal axis that fixes infinity.  Choose the
correspondence in the given longitudinal order and match the two
complementary sides to preserve orientation.  The extension fixes
infinity and hence restricts to an orientation-preserving plane
homeomorphism $\Psi$ taking the ordered axis to this line.  Prescribe
its longitudinal parametrization in the marked order, with $p$
at $u=0$ and all of $\Delta$ strictly within $|u|<1$.  Assign
$u=-16$ and $u=16$ to the first extensions of $\Delta$ in the
two directions that reach radius $2L$.  The central axis segment
through these parameters lies in $\overline B(0,2L)$; by
continuity, for some $d>0$,
$\Psi([-16-d,16+d]\times\{0\})\subset B(0,3L)$.
This compact segment has positive distance from $E$.  Its parameter
interval contains the chart interval $[-12,12]$, the later cuts
at $u=\pm11$, and their longitudinal buffers.
The rest of the line follows the remaining parts of $K$ and,
beyond its endpoints, the prescribed outward rays to infinity.
Outside the marked middle interval the line misses
$\overline B(0,4b_2)$.

For sufficiently small $\tau>0$, the chart
\[
             \Phi(u,v)=\Psi(u,\tau v),\qquad
             (u,v)\in[-12,12]\times[-3,3],
\]
lies in $B(0,4L)$ and misses $E$; its section at zero is strictly inside
$B(0,a_1)$.  We describe the state and obstacle disks to ensure the
prescribed exact slab intersections.  In straight coordinates take
nested polygonal Jordan neighborhoods approximating the preimages of
the disks of radii $2b_2$ and $3b_2$, inside the preimage of
$B(0,4b_2)$.  They contain the preimage of
$\overline B(0,b_2)$ with a margin and avoid a thin strip about the
axis for $|u|\ge2$.  Adjoin thin horizontal rectangles with longitudinal
intervals $[-3,3]$ and $[-8,8]$, respectively, and fill bounded
complementary components.  The polygons and rectangle widths can be
chosen strictly nested and with transverse intersections of boundaries.
Their unions have connected interiors and disjoint simple polygonal
boundary circuits; filling the holes leaves a single outer boundary.
They are therefore nested Jordan disks.

Choose $3\tau$ smaller than these rectangle half-widths.  Beyond the
specified vertical cuts the remaining thin axis strip escapes to
infinity without meeting either union.  Hole filling consequently
does not change its intersections with this strip.  Their images $O,H$
have precisely the slab intersections $|u|\le3$ and $|u|\le8$ required
in Section~\ref{sec:plugs}.  They miss $E$: before filling they miss
this connected set, which reaches the outer circle, and hence $E$ lies
in the unbounded complementary component.  In physical coordinates
the unions lie in $B(0,6L)$, and their filled hulls remain there because
the exterior of that disk is connected.  Fix stubs and a small positive
chart-change tolerance as in Sections~\ref{sec:plugs} and
\ref{sec:affine}.  Decrease the tolerance to preserve all compact
window margins; in particular nominal coordinates in the terminal
windows change by less than $0.1$.  The support of the ambient change
is allowed inside $B(0,50L)$.

It remains to justify exact fitting, rather than only close tracking.
Consider simple profiles converging to the nominal profile in the
marked topology.  In $\Psi$-coordinates take the subarc from its last
hit of $u=-11$ before the core mark to its first hit of $u=11$ after
the mark.  For close enough profiles this is a proper left-to-right
crosscut of a fixed thin rectangle.  Its interior lies between the
vertical sides, and it stays away from the horizontal sides.
Fr\'echet convergence controls its order and backtracking, so these
crosscuts converge in parametrized boundary distance to the axis
crosscut.

There are homeomorphisms of this rectangle converging uniformly to the
identity that take its midline, in order, to the new crosscuts.  Here
is the extension argument.  Split the rectangle along the old midline
and the new crosscut.  On the boundaries of the two old half-rectangles
prescribe correspondences to the two new halves, using the matching
crosscut parametrizations, the identity on the horizontal sides, and
piecewise linear increasing height maps on the vertical sides.  These
maps agree on the common cut and converge uniformly to the identity.
Each new half is a Jordan domain, and its parametrized boundary
converges uniformly to that of the old half.  Choose a fixed interior
point of each old half, which belongs to every sufficiently close new
half, and normalize the conformal maps at that point with positive
derivative.  Rad\'o's uniform convergence theorem for Jordan domains
with Fr\'echet-convergent boundaries
\cite[\S\S1--2, pp.~182--183]{Rado1923} gives convergence of these
maps on the closed disk.  Write them as $F_n,F$, and write $h_n$ for
the prescribed boundary correspondence.  The circle homeomorphisms
$\chi_n=F_n^{-1}h_nF$ converge uniformly to the identity: for
$\zeta\in\partial\mathbb D$,
\[
 |F(\chi_n(\zeta))-F(\zeta)|
       \le\|F_n-F\|_\infty+\|h_n-\mathrm{id}\|_\infty,
\]
and the inverse of the fixed closed Jordan chart $F$ is uniformly
continuous.  Extend these circle homeomorphisms radially by
$re^{\ii\theta}\mapsto r\chi(e^{\ii\theta})$; these are disk
homeomorphisms with the same uniform convergence.  Conjugation by the
conformal charts gives the desired half-rectangle homeomorphisms,
with the prescribed exact values on the common cut and uniform
convergence to the identity.  The cut parametrizations can match all
three interior marks: interpolate the small changes of their terminal
parameters by increasing piecewise linear maps fixing those marks.
Glue them along the cut.  Interpolate the small vertical-side height
changes to the identity in thin exterior collars, keeping the top and
bottom fixed.  The result extends by the identity to a plane
homeomorphism.  After conjugating by $\Psi$, its support can be kept in
$B(0,50L)$ and its displacement tends uniformly to zero.

The portions of the ordered arc omitted by this crosscut correspond,
for close enough profiles, to nominal axis positions with
$|u|>10.5$.  The original compact margins separate them from the changed
obstacle and the central track through $|u|\le9.5$.  The unwanted set
$E$ is separated there as well.  Thus every sufficiently close simple
profile admits the required change within any prescribed positive
tolerance.  First obtain this conclusion on a neighborhood, and then
choose a smaller neighborhood whose closure among simple profiles is
contained in it.  This gives the stated closed-limit property and the
strict prelimit tracking requirements.

Cover the compact family of limiting profiles by finitely many such
open neighborhoods.  If no finite initial constraint list forced
membership in their union, one could choose a profile outside the
union satisfying the first $j$ constraints for every $j$.
Lemma~\ref{lemma:profile-compactness} would give a subsequential limit
inside the union, contradicting openness.  Thus one finite list
suffices.  The construction uses only normalized profile geometry and
the finite trial data, so the same finite menu works at every
sufficiently small $h$.
\end{proof}

On fine retained mesh a fitting traversal supplies vertices near
$u=\pm1.7$ joined along that rim within
$\Phi([-2,2]\times I_B)$, and supplies both side passages in the
nominal windows.  The transverse and backtracking margins, rather than
a smoothness or transversality assumption on the rim, give these
conclusions.  All visits near the second gate annulus belong to
$\Delta$.

\subsection{Local preparations, ranks, and propensities}

For each of the finitely many chart and color choices at a grid center,
apply a separate preparation to the original points and colors.  In
the black case it acts only in a narrow substrip of
$\Phi([-7.5,7.5]\times I_W)$.  It forces a retained white walk from
about $u=-7.1$ to $u=7.1$ with longitudinal backtracking less than $0.2$.
Choose a buffered chain of crossing rectangles along the substrip,
using overlapping progress steps whose chart-longitudinal widths are
much smaller than $0.2$.  Crossings in both directions in the overlaps
join the successive walks.  They can be concatenated in order, with
the desired width and progress margins, through both side-test
terminals.  The corridor argument of Section~\ref{sec:plugs} and the
quenched rectangle bounds give, on a local geometry event with
probability tending to one, a conditional success probability at least
$p_*>0$.  Since there are finitely many normalized models, take a
common positive lower bound.  The good event can be the retained-mesh
event and this lower-propensity event itself.

The preparation uses the finite ordered-coupling principle; see
\cite[Theorem~4.1 and its finite-network discussion]{Edwards1978}.
For completeness, the finite-spin construction needed here is as follows.  If $A$ is increasing under the white order for a finite
product measure $\pi$, positive association implies
$\pi(B\mid A)\ge\pi(B)$ for every increasing set $B$.
There is therefore a coupling $(\omega,\omega')$ with laws
$\pi,\pi(\cdot\mid A)$ and $\omega\le\omega'$.
One may obtain it by a finite flow network: supply $\pi(\omega)$ at
each initial state, demand $\pi(\omega'\mid A)$ at each final state,
and allow an edge exactly when $\omega\le\omega'$.  The inequalities
for increasing sets are precisely the cut inequalities, and the
max-flow criterion supplies the coupling.  Its conditional transition
law is a kernel that only recolors existing black sites white,
with the point geometry fixed.  Fix a deterministic choice of
flow for each finite ordered truth table and probabilities; this also
ensures measurability when the point configuration varies.  Each
bounded local region contains finitely many sites almost surely.

Activate this kernel solely according to its local geometry event,
using an independent seed, and otherwise do nothing and flag the
location.  In particular its activation does not condition on a
fitting rim.  Given the geometry its output law is bounded by
$\max(1,p_*^{-1})$ times ordinary spinning.  Writing
$\sigma^{\mathrm{prep}}$ for the full prepared coloring, its joint
law with the unchanged point configuration consequently satisfies
\[
 \mathcal L(\eta,\sigma^{\mathrm{prep}})
                 \le \max(1,p_*^{-1})\nu,
\]
as measures, uniformly in $h$.
Different candidate outputs need not have a jointly dominated law;
we use the bound for each individual preparation, and for two
disjoint preparations with independent seeds.  The flipped construction
recolors existing white sites black in the opposing strip for white rims,
again fixing the point geometry.

At a fit the preparation avoids the rim.  It keeps both gate circuits
with at most two black vertices and supplies a single white traversal
whose root portion and side portions are compatible.  Together with
the original black traversal, some pair of root ranks therefore has
both side successes.  The ranks can be truncated uniformly with
arbitrarily small error.  Indeed distinct crossing components of a
root window supply vertex-disjoint same-color paths across its middle
slab.  On good geometry an opposing transverse bar has quenched
probability at least $q>0$.  Thus the probability of one such
monochromatic crossing is at most $1-q$, and the conditional BK
inequality bounds the probability of $k$ disjoint crossings by
$(1-q)^k$.  Transfer this bound to the prepared output using its
density bound, and discard the negligible geometry failure.  A fixed
large rank cutoff works for the finite collection of windows.

We next discard actual successes whose ordinary conditional
propensities are too small.  For any event $A$ and sigma field
$\mathcal F$,
\begin{equation}\label{pin:small-propensity}
 \nu\bigl(A\cap\{\nu(A\mid\mathcal F)<t\}\bigr)
 =\EE_\nu\bigl[\ind_{\{\nu(A\mid\mathcal F)<t\}}
                 \nu(A\mid\mathcal F)\bigr]\le t.
\end{equation}
Apply this first to $L^+$ and $L^-$ conditional on the inner state,
and then to each white side success conditional also on the local
exterior point geometry.  Transfer each estimate by the preparation
density and take a finite union over the retained ranks and models.
Choosing $\beta>0$ small makes the total cost arbitrarily small.
Finally, for each resulting type discard its actual prepared
occurrence of $G$ when $\nu(G)<\beta'$.  Its probability is at most
$C\beta'$ per type.  Choose $\beta'>0$ after the finite rank
truncation.  The remaining prepared states have exactly the
qualifications in Definition~\ref{pin:qualification}.

All discard flags are local.  A side-test propensity conditional on the
inner state is a measurable function of that state and fixed local
test; its quenched counterpart additionally uses only the exterior
point restriction needed by its retained paths.  The preparations,
rank counts, actual side successes, and admission indicators use no
hypothetical long connector.  Their range is bounded by $CL$ in
normalized units.  Stubs are not queried by these flags.

\subsection{From rare flags to dense pins}

\begin{proposition}[Density of qualified pins]\label{prop:dense-pins}
There exist a finite normalized menu and fixed $\beta,\beta'>0$ such
that the following holds in every bulk patch with
$\overline B(x,6R)$ available.  With probability tending to one in the
ordered limit $\eps\to0$, then $h\to0$, every subarc of diameter at
least $r_h$ in $\overline B(x,2R)$ on a rim reaching outside
$B(x,4R)$ contains a qualified pin.  The statement holds simultaneously
for all such subarcs and both colors.  It also holds for the coupled
bulk settings on their patch-agreement event.
\end{proposition}

\begin{proof}
For either color, the preceding local tests have the following
deterministic implication on patch agreement and good mesh: every
eligible marked rim vertex within distance one of a grid center fits
some qualified choice unless that center is flagged.  The trial part
of its limiting flag probability is \eqref{pin:trials}.  For fixed $L$
all other costs can be made arbitrarily small.  The order of choices is
essential.  Choose $M,n$ first.  Allocate a summable inverse-test budget,
which determines the finite profile menu and its finite constraint
list.  For that menu fix the preparation corridors and their positive
lower bounds, then truncate ranks, and finally choose the two
propensity thresholds.  Every choice is in normalized units and
independent of $h$.  Admission itself need not stabilize here.

Tile the plane by blocks of side $Lh$.  A block is bad if its enlargement
by distance $2Lh$ contains at least $c_0L$ flagged grid centers, where
$c_0>0$ is a sufficiently small absolute constant.  There are $O(L^2)$
candidate centers in that enlargement, so the first-moment bound gives
\begin{equation}\label{pin:flag}
 \limsup_{\eps\to0}\PP(\text{a specified block is bad})
 \le CL\bigl[Cn^2(CM^{-1-c})^{n-2}+\zeta\bigr],
\end{equation}
where $\zeta$ is the sum of the adjustable additional costs.  Since
$L=30M(10^4M)^{n-1}$, first taking $M$ large and then $n$ large makes
the contribution $L n^2(CM^{-1-c})^{n-2}$ as small as desired.
Choose $\zeta$ afterward.  The dependence range of these block events
is an absolute constant in block units, including the independently
assigned preparation seeds.

An eligible subarc of diameter at least $r_h$ without a qualified pin
forces a connected chain of bad blocks of length at least
$c r_h/(Lh)$.  Choose two points of the subarc at distance at least
$r_h/2$ and follow the portion between them.  Since $r_h/(Lh)\to\infty$,
at least one of these endpoints lies outside the $2Lh$ enlargement of
each visited block.  From a visit to the block, follow the subarc toward
that endpoint until it first leaves the enlargement.  Even if the visit
only clips a corner of the block, this portion travels distance at least
$2Lh$.  On fine mesh, its initial passage over distance $Lh$ supplies
vertices within distance $h$ of at least $c_0L$ distinct grid centers,
all still inside the enlargement.  Every such center is flagged by the
deterministic implication above.  Thus every visited block is bad.
The visited blocks connect the two chosen endpoints; erasing repetitions
gives a chain of distinct neighboring bad blocks, whose length is at
least $c r_h/(Lh)$ by their separation.

There are exponentially many block chains of a specified length, and
only $O(h^{-2})$ possible starting blocks in the fixed patch.  From any
chain of length $k$ choose at least $k/C_0$ blocks whose flag-data
regions are pairwise disjoint.  Their bad events are independent under
the full-plane marked Poisson law with the local seeds.  If the bound
in \eqref{pin:flag} is sufficiently small, exponential path counting
therefore gives an upper bound of the form
\[
          C_Rh^{-2}\exp\{-c r_h/(Lh)\}\longrightarrow0.
\]
At each fixed $h$ the chain length and all local scale lists are finite;
take the microscopic limit before this estimate.  Discarding the
vanishing patch-agreement and mesh failures completes the proof.
\end{proof}
